\subsection{Ablation y robustez}
Para estimar qué información aporta cada grupo de variables se entrenó un mismo HistGradientBoosting con la configuración seleccionada por integrante 2 (180 iteraciones, 15 hojas, tasa de aprendizaje 0.1, mínimo de 80 ejemplos por hoja, regularización L2 de 5.0) con cuatro subconjuntos de las 33 características: A0, la tasa actual (\texttt{median\_fee\_rate}, \texttt{umbral\_local} y \texttt{fee\_relativa}); A, las variables de comisiones (16); B, comisiones y actividad de la red, incluidas las comisiones totales como se planteó en H2 (29); y C, el conjunto completo con variables de calendario (33). Se emplearon los mismos folds temporales, el mismo \textit{gap} y el mismo criterio de selección del umbral que en los modelos de ensamble. La Tabla~\ref{tab:ablation} resume los resultados en validación.

\begin{table}[htbp]
\caption{Ablation en validación cronológica (62\,808 ejemplos). La persistencia solo emite etiquetas, por lo que su AP no es comparable.}
\label{tab:ablation}
\centering
\small
\begin{tabular}{lcccc}
\hline
Conjunto & Var. & AP & F1 & ROC-AUC \\
\hline
A0: tasa actual & 3 & 0.5345 & 0.5389 & 0.7492 \\
A: comisiones & 16 & 0.6224 & 0.5692 & 0.8178 \\
B: + actividad & 29 & 0.6291 & 0.5653 & 0.8202 \\
C: + calendario & 33 & 0.6298 & 0.5713 & 0.8217 \\
Persistencia & -- & -- & 0.5400 & -- \\
\hline
\end{tabular}
\end{table}

Las diferencias se evaluaron con un \textit{bootstrap} emparejado por bloques móviles de 1008 bloques sobre validación (1000 remuestreos), para respetar la dependencia entre bloques consecutivos. Incorporar la historia de las comisiones (A frente a A0) aumentó la AP en 0.088 (IC 95\,\%: [0.075, 0.099]). Añadir la actividad de la red (B frente a A) la aumentó en 0.007 ([0.003, 0.011]). Las variables de calendario no cambiaron la AP (C frente a B: +0.001, [$-0.005$, 0.005]). En F1 las diferencias entre A, B y C fueron pequeñas y de signo variable, porque dependen del umbral elegido en cada caso. El mismo orden se observó en los tres folds de entrenamiento.

El rendimiento se analizó también por trimestre y según la volatilidad reciente, medida como el coeficiente de variación de las últimas 144 tasas, con terciles fijados en entrenamiento. Entre trimestres, el F1 osciló entre 0.50 y 0.64 de forma paralela en todos los métodos, siguiendo la prevalencia de la clase alta, y la ventaja sobre la persistencia fue positiva en todos ellos. En el tercil de mayor volatilidad, el F1 del modelo C bajó a 0.554, frente a 0.567 y 0.582 en los terciles bajo y medio, y la AP a 0.596, frente a 0.626 y 0.643.

En el análisis de errores, el 19.8\,\% de los bloques de validación correspondió a un cambio de clase, y los episodios de comisión alta duraron en promedio 2.17 bloques. La persistencia falla, por construcción, en todas esas transiciones. El modelo C anticipó el 35\,\% de los finales de episodio, pero solo el 19\,\% de los inicios, a costa de errar en el 4.9\,\% de los tramos bajos estables y en el 8.5\,\% de los altos estables. La tasa de error fue del 8\,\% cuando la tasa del bloque siguiente quedó por debajo de la mitad del umbral y del 72\,\% cuando lo superó por menos del 5\,\%.
